\section{Damped removal of smoothing and higher-moment stopping}
\label{sec:damped-unsmoothing}
The first-order comparison in
Lemma~\ref{lem:fixed-count-damped-comparison} takes the absolute value
of the entire change of observable before using spectral damping.
We instead retain a finite Taylor polynomial of that change inside
the chronological collision word. Every term of this polynomial
still has damping over the full deterministic collision interval.
Only its fourth-order remainder is estimated without damping.
A second, finer smoothing scale is used for the finitely many
insertions; it does not shrink the analytic domain of the twisted
collision powers. We also improve the actual stopping estimate by
using a fixed higher moment of the visit count.

All constants below are uniform in $R$. Their dependence on a fixed
moment order is allowed. We neither let that order grow with $n$
nor use an all-orders analytic expansion. Write
$\|u\|_{\mathcal V}=\|u\|_\infty+\|u\|_{\mathrm{BV}}$ and
$L_\delta=1+|\log\delta|$. Expectations without a subscript are with
respect to the stationary collision probability $\nu$.

\subsection{Cumulants of a small bounded-variation residual}
\begin{lemma}[Small-mass connected correlations]
\label{lem:small-mass-cumulants}
Fix $q\ge2$ and $B<\infty$. If $u$ is real,
$\|u\|_{\mathcal V}\le B$, and $\|u\|_1\le B\delta$ for
$0<\delta<1/2$, then
\begin{align}
 &\sum_{\mathbf j\in\Z^{q-1}}
 \left|\operatorname{cum}(u,u\circ T_R^{j_2},\ldots,
                                      u\circ T_R^{j_q})\right|
       \le C_{q,B}\delta L_\delta^{q-1},
       \label{eq:small-mass-connected}\\
 &|\operatorname{cum}_q(S_{m,R}u)|
       \le C_{q,B}m\delta L_\delta^{q-1}\qquad(m\ge1).
       \label{eq:small-mass-cumulant-length}
\end{align}
If in addition $\int u\dd\nu=0$, then
\begin{equation}\label{eq:small-residual-fourth}
 \mathbb E|S_{m,R}u|^4
       \le C_B\bigl(m^2\delta^2L_\delta^2
                          +m\delta L_\delta^3\bigr).
\end{equation}
The same fourth-moment bound holds on every deterministic integer
interval, also with its orientation reversed.
\end{lemma}
\begin{proof}
In each term of the partition formula
\eqref{eq:finite-order-cumulant-definition}, the block containing the
first factor has moment at most $B^{q-1}\|u\|_1$ in absolute value;
all other block moments are bounded. Thus the joint cumulant is at
most $C_{q,B}\delta$, independently of the times. The independent-block
argument in Proposition~\ref{prop:connected-collision-summability}
also bounds it by $C_{q,B}e^{-\gamma D}$, where
$D=\operatorname{diam}\{0,j_2,\ldots,j_q\}$ and $\gamma>0$
depends on the fixed order.

There are at most $(2J+1)^{q-1}$ anchored tuples with $D\le J$.
Choose $J$ to be a sufficiently large fixed multiple of $L_\delta$.
Inside this set use the bound $C\delta$; outside it sum the
exponential bound, absorbing the fixed-degree lattice counting
polynomial in half of the exponential. This gives
$C_{q,B}\delta L_\delta^{q-1}$. The exact anchored multiplicity
$(m-D)_+$ is at most $m$, proving
\eqref{eq:small-mass-cumulant-length} without a boundary error
independent of $\delta$.

For a centered real variable $X$, the finite moment--cumulant identity
is $\mathbb E X^4=\operatorname{cum}_4(X)+
3\operatorname{cum}_2(X)^2$. It follows either by expanding the
exponential of its cumulant generating polynomial through degree
four, or directly by the partition formula. Apply
\eqref{eq:small-mass-cumulant-length} at orders two and four.
Stationarity translates an arbitrary integer interval, and reversal
does not change its sum.
\end{proof}

\begin{lemma}[Every fixed even collision moment and its maximum]
\label{lem:fixed-even-maximal}
Fix an integer $p\ge2$. For a real centered $u$ with
$\|u\|_{\mathcal V}\le B$ and $m\ge1$,
\begin{equation}\label{eq:fixed-even-maximal}
 \mathbb E\left|\sum_{j=r}^{r+m-1}u\circ T_R^j\right|^{2p}
 +\mathbb E\max_{0\le l\le m}
       \left|\sum_{j=r}^{r+l-1}u\circ T_R^j\right|^{2p}
       \le C_{p,B}m^p.
\end{equation}
The location $r\in\Z$ is arbitrary. The estimate holds in reverse
order, for the vector $h_R$, and under a nonnegative initial density
bounded by a fixed constant.
\end{lemma}
\begin{proof}
By Proposition~\ref{prop:connected-collision-summability}, every
cumulant of $S_m u$ of order $2\le j\le2p$ is $O_{p,B}(m)$.
The coefficient identity
\[
 \mathbb E X^{2p}=\sum_{\pi\in\mathcal P_{2p}}
                         \prod_{B\in\pi}\operatorname{cum}_{|B|}(X)
\]
is obtained by exponentiating the generating polynomial through that
fixed degree. For centered $X$, partitions with singleton blocks
vanish. Every remaining partition has at most $p$ blocks. The
moment of $S_m u$ is consequently at most $C_{p,B}m^p$.
No convergence of an infinite cumulant series is used.

Enlarge $m$ to a power of two $M<2m$. Every prefix is a union of at
most one aligned dyadic interval at each scale. At interval length
$l=M2^{-j}$ the $2p$th moment of the maximum interval sum is at most
$C_{p,B}(M/l)l^p$. Minkowski's inequality over scales therefore gives
\[
 \left\|\max_{0\le l\le M}|S_lu|\right\|_{2p}
 \le C_{p,B}\sum_{j=0}^{\log_2 M}
                 (M^p2^{-j(p-1)})^{1/(2p)}
 \le C_{p,B}\sqrt M.
\]
The geometric series converges because $p>1$. Translation and
reversal use invariance, not a changed initial law. A bounded initial
density is handled by domination. Summing the four component bounds
proves the vector assertion.
\end{proof}

\subsection{Higher-moment windows for the genuine return clock}
Retain $N^\pm_{j,R}$, $m_j=\lfloor j/c_*\rfloor$ and
$\mathsf H_{r,s,R}$ from Section~\ref{sec:marked-return-band}.
As there, the zero-return clocks and empty sums are zero.

\begin{proposition}[Higher-moment marked stopping]
\label{prop:higher-moment-marked-stopping}
For each fixed integer $p\ge2$, all $j\ge0$ and $b\ge2$,
\begin{equation}\label{eq:higher-moment-clock}
 \nu_R^*\{|N^\pm_{j,R}-m_j|>b\}
             \le C_p\frac{(j+b)^p}{b^{2p}}.
\end{equation}
Put $s_p=p/(4p+1)$. For every $n\ge2$, $0\le k\le n$ and
$v\in\R^4$, the actual marked transform obeys
\begin{align}
 &\left|C^{[k],a}_{n,R}(v)
 -\int_M a(y)e^{iv\cdot\mathsf H_{m_k,m_{n-k},R}(y)/\sqrt n}
                       \dd\nu(y)\right|\notag\\
 &\hspace{18mm}\le C_pM_a(1+|v|)n^{-s_p}.
 \label{eq:higher-moment-marked-stopping}
\end{align}
No bound on $V_a$ is needed for this stopping comparison.
\end{proposition}
\begin{proof}
The centered visit count
$\sum_{i=1}^L(\eta_R-c_*)\circ T_R^{\pm i}$ has $2p$th moment at
most $C_pL^p$ under $\nu_R^*$, by
Lemma~\ref{lem:fixed-even-maximal} and $\nu_R^*\le c_*^{-1}\nu$.
The exact visit equivalences in
Lemma~\ref{lem:two-sided-clock-window}, followed by Markov's
inequality at integer thresholds within a bounded distance of
$m_j\pm b$, give \eqref{eq:higher-moment-clock}. The required
deviation is at least $c_*b-O(1)$ and $L\le C(j+b)$.
A negative lower threshold yields an empty event; $j=0$ is immediate.
Enlarge the constant for bounded $b$.

Recenter at the actual marked return using
\eqref{eq:marked-recentering}. For $2\le b\le n$, the union of the
two clock exceptions has probability at most $C_pn^p/b^{2p}$.
Off this union, the difference from the deterministic interval is
bounded by maxima on two fixed collision windows of length $O(b)$.
Lemma~\ref{lem:fixed-even-maximal} bounds their second moments by
$C_pb$, including windows crossing collision time zero. On the
exceptional set use only the bound two for a difference of phases.
Consequently the error is at most
\[
 C_pM_a\left(\frac{n^p}{b^{2p}}+|v|\sqrt{\frac bn}\right).
\]
Taking $b=\lceil n^{(2p+1)/(4p+1)}\rceil$ gives
\eqref{eq:higher-moment-marked-stopping}. Finitely many small $n$
are absorbed in the constant. The two clocks are not assumed
independent, and the weight is never evaluated at a deterministic
surrogate for the marked return.
\end{proof}

\subsection{A cubic correction retained inside the damped word}
Let $r,s\ge0$, $m=r+s\ge1$, and define
\[
 I_{r,s,R}^a(z)=\int_M a(y)
                         e^{iz\cdot\mathsf H_{r,s,R}(y)}\dd\nu(y).
\]
The insertion need not be positive. The next result uses the coarse
scale $\delta$ only for $h_{R,\delta}$ and the fine scale $\epsilon$
only for the insertion and a finite number of residual factors.

\begin{theorem}[Damped finite-order removal of smoothing]
\label{thm:damped-cubic-unsmoothing}
Fix $Q\ge3$, and suppose $\delta\le\delta_0$ and real $z$ satisfy
\eqref{eq:finite-jet-damping-domain}. For every $0<\epsilon<1/2$,
\begin{align}
 |I_{r,s,R}^a(z)|\le{}&
 CM_a\epsilon^{-8}(1+m|z|)^3e^{-cm|z|^2}+C\epsilon V_a\notag\\
 &+CM_a\epsilon(1+m|z|)^3\notag\\
 &+CM_a|z|^4\bigl(m^2\delta^2L_\delta^2
                              +m\delta L_\delta^3\bigr).
 \label{eq:damped-cubic-unsmoothing}
\end{align}
The constants are independent of the relative location of the mark
and of the fine scale. No analytic perturbation in the fine-scale
observable is asserted.
\end{theorem}
\begin{proof}
The assertion at $z=0$ follows by increasing $C$, so take $z\ne0$.
Put $\xi=z/|z|$, $u=\xi\cdot(h_R-h_{R,\delta})$,
$a_\epsilon=\mathsf S_\epsilon a$ and
$u_\epsilon=\mathsf S_\epsilon u$. Mean-preserving smoothing gives
\[
 \int u\dd\nu=0,\quad \|u\|_{\mathcal V}\le C,\quad
 \|u\|_1\le C\delta,\quad
 \|u-u_\epsilon\|_1\le C\epsilon.
\]
Also $\|a-a_\epsilon\|_1\le C\epsilon V_a$ and
$\|a_\epsilon\|_\infty\le M_a$. Replace the insertion first,
while the full exponential still has modulus one. This costs exactly
the second term of \eqref{eq:damped-cubic-unsmoothing}.

Translate the interval to collision times $0,\ldots,m-1$. Its mark
is now at time $r\in[0,m]$. With $X=|z|S_m u$, use the real-variable
Taylor formula
\begin{equation}\label{eq:cubic-phase-remainder}
 e^{iX}=\sum_{j=0}^3\frac{(iX)^j}{j!}+\mathcal R_4(X),
 \qquad |\mathcal R_4(X)|\le |X|^4/24.
\end{equation}
Multiplication by the coarse phase
$e^{iz\cdot S_m h_{R,\delta}}$ does not change this bound.
Lemma~\ref{lem:small-mass-cumulants} therefore bounds the remainder
integral by the last term of \eqref{eq:damped-cubic-unsmoothing}.
This estimate is for the original residual $u$, not for a residual
with a presumed independent law.

For $1\le j\le3$, expand $(S_m u)^j$ as a sum over its $m^j$
ordered index tuples, allowing repetitions. Within each product,
replace each factor $u\circ T_R^t$ by $u_\epsilon\circ T_R^t$.
The other factors and the coarse phase are bounded. Telescoping and
invariance bound its integrated error by $C_jM_a\epsilon$.
After multiplying by $|z|^j/j!$ and summing, this costs at most
$CM_a\epsilon(1+m|z|)^3$. No long pullback BV norm is used.

It remains to bound each corrected product. Write
$L_z=\mathcal L_{R,\delta,z}$. Sort its $j$ residual insertion times
together with the mark, retaining repeated times, as
$0\le t_1\le\cdots\le t_d\le m$, where $d=j+1\le4$.
Exactly one $b_i$ is $a_\epsilon$; the others are $u_\epsilon$.
The chronological word is
\begin{equation}\label{eq:corrected-chronological-word}
 \ell\!\left(
 L_z^{m-t_d}M_{b_d}L_z^{t_d-t_{d-1}}M_{b_{d-1}}
 \cdots L_z^{t_2-t_1}M_{b_1}L_z^{t_1}\nu\right).
\end{equation}
For $d=1$ this means
$\ell(L_z^{m-t_1}M_{b_1}L_z^{t_1}\nu)$.
The nonnegative power lengths sum to $m$. Smooth multiplication costs
at most $CM_a\epsilon^{-2}$ for the mark and $C\epsilon^{-2}$
for each residual. Corollary~\ref{cor:finite-jet-damping} thus bounds
\eqref{eq:corrected-chronological-word} by
$C_jM_a\epsilon^{-2(j+1)}e^{-cm|z|^2}$. Summing the finitely
many degrees and all their index tuples proves the first term of
\eqref{eq:damped-cubic-unsmoothing}.

A repeated time is a zero-length power, and a mark at $0$ or $m$
is an endpoint multiplier; both are covered by the same power bound.
There is never a perturbation of $L_z$ by $u_\epsilon$. In particular,
the fine scale does not enter the damping-domain restrictions.
\end{proof}

\paragraph{Why the new estimate has a different scale balance.}
The undamped observable loss in the earlier comparison was proportional
to $|z|\sqrt{m\delta L_\delta}$. In
\eqref{eq:damped-cubic-unsmoothing}, the lower-degree changes remain
inside a word damped for all $m$ collisions; its undamped residual
begins with $|z|^4m^2\delta^2L_\delta^2$. The polynomial cost of
the much finer insertions is harmless on an annulus with a growing
rescaled inner radius. This argument concerns the characteristic
function itself at a deterministic collision length. The passage to a
specified actual return count uses
\eqref{eq:higher-moment-marked-stopping}, not an average over counts.
